% Rascunho para a Etapa 6: secao retirada da Etapa 5 (resultados com orcamento por tempo).
% Atualizar com a rodada de 50.000 cronogramas e o Gurobi no J60.
% Referencias usadas: \cite{rko} Chaves et al. (2025); \cite{brkga} Goncalves e Resende (2011).

\section*{Comparação preliminar com as metaheurísticas}

Para situar o método exato e responder à questão levantada em sala sobre o desempenho de
algoritmos genéticos neste problema, as mesmas instâncias foram resolvidas com três metaheurísticas
do framework RKO (\emph{Random-Key Optimizer}) \cite{rko}: o algoritmo genético por chaves
aleatórias (GA), o BRKGA \cite{brkga} e o RKO completo, que executa nove metaheurísticas em
paralelo. Os três usam o mesmo decodificador (chaves $\to$ lista de prioridade $\to$ SSGS, com ou
sem FBI), com 10 execuções por instância (10~s no J30; 30~s em cinco instâncias J60 com $RS = 0{,}2$),
na mesma máquina e com parada ao atingir o ótimo.

\begin{table}[ht]
\centering
\small
\caption{Metaheurísticas: RPD médio $= 100\,(\cmax - C^*)/C^*$ e execuções que atingiram $C^*$.}
\label{tab:mh}
\begin{tabular}{llrrr}
\toprule
Decodificador & Método & RPD J30 (\%) & RPD J60 (\%) & Execuções no ótimo \\
\midrule
SSGS + FBI & GA    & 0,29 & 0,92 & 94/150 \\
           & BRKGA & 0,24 & 1,00 & 95/150 \\
           & RKO   & 0,28 & 1,00 & 94/150 \\
\midrule
Só SSGS    & GA    & 0,43 & 2,00 & 79/150 \\
           & BRKGA & 0,51 & 2,07 & 70/150 \\
           & RKO   & \textbf{0,38} & \textbf{1,69} & \textbf{86/150} \\
\bottomrule
\end{tabular}
\end{table}

\vspace{-0.4em}
\textbf{Métodos de busca.} Com o FBI, os três métodos são estatisticamente equivalentes
(Friedman, $p = 0{,}81$): quando o GA usa a mesma representação e o mesmo decodificador, não se
observa a dificuldade esperada. Sem o FBI, o RKO obtém os melhores valores, mas a diferença não é
significativa após a correção para comparações múltiplas (Wilcoxon com Holm, $p = 0{,}13$). O fator
de maior efeito é o decodificador: no J60, o FBI reduz o RPD médio de 1,7--2,1\% para 0,9--1,0\% nos três
métodos ($p < 0{,}01$). Uma limitação é o orçamento: em 10--30~s cada execução avalia milhões de
cronogramas, muito acima dos 5.000 a 50.000 usados na literatura, o que tende a aproximar os métodos.

\textbf{Exato $\times$ metaheurística.} No J30, a melhor solução do RKO é igual ou melhor que a do
Gurobi nas dez instâncias, obtida em 0,58~s em média, contra 259~s do Gurobi. Em j3013\_1, o RKO
alcançou o ótimo (58) em 7 das 10 execuções, enquanto o Gurobi terminou em 59 após 600~s; em
j3029\_8, nenhum dos dois atingiu o ótimo (80). A vantagem do método exato está na prova de
otimalidade, obtida em 7 das 10 instâncias.

